% =====================================================================
%  Resumo (LNCS: 150 a 250 palavras) e palavras-chave
% =====================================================================
\begin{abstract}
As Rotinas da Amazon Alexa permitem associar um evento a uma sequência de
ações executadas pelo assistente e pelos dispositivos conectados. Elas são
criadas por uma pessoa, no aplicativo ligado à sua conta, e executadas em
espaços compartilhados da casa, onde afetam todos que estão presentes. Este
artigo avalia a metacomunicação do aplicativo Amazon Alexa na criação e no
gerenciamento de rotinas, considerando a casa como usuário coletivo. Para
isso, aplicamos o Método de Inspeção Semiótica com uma lente para múltiplos
interlocutores, que distingue o que a interface comunica a quem cria a
rotina e o que comunica sobre quem é afetado por ela. Os achados são
triangulados com avaliações públicas do aplicativo nas lojas brasileiras.
\pendente{resumir os principais resultados e implicações de design após a
inspeção}
\keywords{Engenharia Semiótica \and Método de Inspeção Semiótica \and
Assistentes de voz \and Casas inteligentes \and Rotinas \and IHC}
\end{abstract}
